\section{Order, generated as a partial order on charges}
\begin{proposition}[The knowledge order $\le_k$]\label{thm:kle}\lean{Obsmetrologia.Reading.KLE}\leanok\uses{thm:four-bilattice}
\texttt{KLE a b} is the proposition that the decidable knowledge comparison
$\texttt{leK}\,a\,b$ holds, i.e.\ that charge $a$ carries no more knowledge than
$b$. It is declared reducible so the underlying boolean is visible to
\texttt{decide}. This is the relation settling generates and climbs.
\end{proposition}
\begin{theorem}[$\le_k$ is reflexive]\label{thm:kle-refl}\lean{Obsmetrologia.Reading.kle_refl}\leanok\uses{thm:kle}
Every charge stands in the knowledge order to itself: $\texttt{KLE}\,a\,a$.
\end{theorem}
\begin{theorem}[$\le_k$ is transitive]\label{thm:kle-trans}\lean{Obsmetrologia.Reading.kle_trans}\leanok\uses{thm:kle}
If $a \le_k b$ and $b \le_k c$ then $a \le_k c$; proven by exhausting the
finite charge lattice with \texttt{decide}.
\end{theorem}
\begin{theorem}[$\le_k$ is antisymmetric]\label{thm:kle-antisymm}\lean{Obsmetrologia.Reading.kle_antisymm}\leanok\uses{thm:kle}
If $a \le_k b$ and $b \le_k a$ then $a = b$. Together with reflexivity and
transitivity this makes $\le_k$ a genuine \emph{partial order}, not merely a
preorder --- the antisymmetry is exactly the thinness that later forces
lax-idempotence to collapse to strict idempotence.
\end{theorem}
